% sections/intro.tex -- Section 1 (paper_v2_vn, translated in L4 from paper_v2's
% sections/intro.tex; terms as in glossary_vn.tex). The claims stay within paper_v2's; every number
% is \val{} from numbers.tex.
\section{Giới thiệu}
\label{sec:intro}

Các hệ thống gợi ý thường biết nhiều hơn việc ai đã mua gì: sản phẩm thuộc về các danh mục, khách
hàng mua kèm một số sản phẩm với nhau và người dùng nhận xét về chúng. Một đồ thị không đồng nhất
chứa được mọi quan hệ này và meta-path là một cách phổ biến để khai thác nó:
\code{PUP}, sản phẩm--người dùng--sản phẩm, nối hai sản phẩm được cùng một người dùng nhận xét.
HAN \citep{wang2019heterogeneous} biến mỗi meta-path thành một đồ thị và gán trọng số cho các
meta-path bằng cơ chế chú ý. Nhiều mô hình gợi ý được xây dựng trên meta-path
\citep{hu2018leveraging,shi2020hgrec,mu2021fhrec,yan2021attention}. Việc chọn meta-path có ảnh
hưởng: hiệu năng của các mô hình gợi ý dựa trên meta-path
``\foreignlanguage{english}{highly sensitive to the meta-paths they use}'' (rất nhạy với các meta-path mà chúng sử dụng) \citep{ning2022automatic}.

Các meta-path thường được chọn thủ công, một việc đòi hỏi kiến thức chuyên gia và không dùng lại
được từ bộ dữ liệu này sang bộ dữ liệu khác \citep[][Phụ lục~D]{han2020genetic,li2024longrange}.
Hướng phát hiện tự động nhằm đảm nhận việc lựa chọn này (Mục~\ref{sec:related}). LMSPS
\citep{li2024longrange}, một phương pháp tìm kiếm được thực hiện trên một không gian lớn các
meta-path ứng viên, cho biết các meta-path nó tìm được nâng Micro-F1 của HAN thêm
\val{lmsps-t4-gain-min} đến \val{lmsps-t4-gain-max} điểm trên bốn bộ dữ liệu cho bài toán phân loại nút
\citep[Bảng~4]{li2024longrange}. Liệu chúng có giúp ích cho HAN trong bài toán gợi ý hay không
thì, theo hiểu biết của chúng tôi, chưa được kiểm chứng và câu trả lời không hiển nhiên, vì ba lý do.
Thứ nhất, LMSPS đánh giá meta-path bằng mạng riêng mà nó dùng khi tìm kiếm chứ không phải bằng HAN,
nên những meta-path tốt nhất với mạng đó chưa chắc là tốt nhất với HAN. Thứ hai, trong tìm kiếm kiến trúc mạng nơ-ron, các phương pháp tìm kiếm thường không tốt
hơn các lựa chọn ngẫu nhiên \citep{li2019random,yu2020evaluating}. Thứ ba, việc tìm kiếm tốn giờ
GPU, trong khi thiết kế thủ công không cần giờ GPU nào.

Vì vậy, chúng tôi đặt câu hỏi: \emph{so với thiết kế thủ công, tìm kiếm meta-path cho gợi ý
sản phẩm dựa trên HAN có đáng với chi phí không?} Trên Amazon-Meta \citep{leskovec2007viral},
chúng tôi so sánh hai cách chọn meta-path cho HAN từ cùng \val{search-cands} ứng viên.
Cách thứ nhất là thủ công, với các tập được thiết kế với sự hỗ trợ
của AI theo các nguyên tắc nêu rõ và cố định trước khi có bất kỳ kết quả nào. Cách thứ hai là tìm
kiếm bằng LMSPS, được điều chỉnh cho bài toán dự đoán liên kết người dùng--sản phẩm và chạy
\val{search-runs} lần cho mỗi kích thước tập, vì một lần chạy không xếp hạng được các tập tốt nhất
của nó. Để biết việc chọn meta-path ảnh hưởng đến đâu, chúng tôi dùng các tập chọn ngẫu nhiên từ
cùng các ứng viên làm đối chứng. Mỗi tập được huấn luyện trên HAN và metapath2vec
\citep{dong2017metapath2vec} và thử nghiệm trên cùng các cạnh giữ lại. Độ đo chính là Hit@10 lấy
mẫu, cố định từ trước, còn xếp hạng toàn bộ dùng để kiểm tra. Mô hình phổ biến và một GAT không
phân biệt kiểu làm mô hình cơ sở.

Trên lược đồ này, câu trả lời tùy thuộc vào bối cảnh, vì kết quả của tìm kiếm khá gần với thiết kế thủ công.
Theo độ đo đã cố định từ trước, các tập tìm được kém các tập thủ công \val{han-4-hit-relabs}\% và
\val{han-6-hit-relabs}\% trong HAN với bốn và sáu meta-path (\val{han-4-hit-pt} và \val{han-6-hit-pt})
và kém \val{mtv-4-hit-relabs}\% và \val{mtv-6-hit-relabs}\% trong metapath2vec. Trong HAN, chúng ngang bằng
khi xếp hạng toàn bộ, sau khoảng \val{search-total-h}~giờ~GPU tìm kiếm. Vì vậy, thiết kế thủ công là lựa chọn
tốt hơn khi có chuyên gia làm được việc này trong thời gian ngắn hơn thời gian tìm kiếm, còn tìm kiếm cho
các tập gần tương đương khi không có chuyên gia. Các tập tìm được thua chủ yếu ở sách và nhạc, những sản phẩm
mà các meta-path danh mục của chúng không đến được. Đối chứng ngẫu nhiên cho thấy trên lược đồ này việc chọn
meta-path ít ảnh hưởng: trong HAN, các tập chọn ngẫu nhiên đạt kết quả ít nhất ngang bằng cả hai cách. Ngoài
ra, GAT vượt cả hai mô hình dùng meta-path. Có hai giới hạn vẫn cần xem xét. Thứ nhất, khi giai đoạn khởi động
(warm-up) của trọng số meta-path trong HAN ngắn hơn, tập tìm được lại vượt tập thủ công ở sáu meta-path, trong
một phép kiểm tra chỉ dùng một seed (Mục~\ref{sec:res-sensitivity}). Thứ hai, do lược đồ nhỏ, bộ dữ liệu này
không kiểm chứng được trường hợp mà tìm kiếm tự động nhắm tới, đó là những lược đồ quá lớn để thiết kế thủ công
hay chọn ngẫu nhiên.

Các đóng góp của chúng tôi gồm:
\begin{itemize}
  \item Một phiên bản Amazon-Meta đã làm sạch cho bài toán dự đoán liên kết, trong đó việc gộp các
        ấn bản trùng lặp giảm tỷ lệ cạnh thử nghiệm có bản sao trong tập huấn luyện từ
        \val{before-testleak-pct}\% xuống \val{after-testleak-pct}\% (Mục~\ref{sec:data}).
  \item LMSPS được điều chỉnh cho bài toán dự đoán liên kết, kèm các biện pháp chống rò rỉ dữ liệu
        và một cách kết hợp nhiều lần chạy tìm kiếm thành một quyết định
        (Mục~\ref{sec:method}).
  \item Một phép so sánh có kiểm soát giữa các tập meta-path tìm được và thủ công trên HAN và
        metapath2vec, với đối chứng ngẫu nhiên, trong đó các tập và độ đo chính được cố định từ
        trước (Mục~\ref{sec:setup} và~\ref{sec:results}).
  \item Một góc nhìn chi phí--hiệu quả về tìm kiếm so với thiết kế thủ công, với bằng chứng ở mức
        meta-path cho thấy các tập tìm được thua ở đâu (Mục~\ref{sec:results}).
\end{itemize}
